\section*{Group actions and invariant generators (original Section1)}
An action (also called a representation) of a group $G\subseteq\GL_n(\C)$ on an $m$-dimensional complex vector space $V$ is a group homomorphism $\pi:G\to\GL_m(\C)$, that is, an association of an invertible $m\times m$ matrix $\pi(g)$ for every group element $g\in G$, satisfying $\pi(g_1g_2)=\pi(g_1)\pi(g_2)$ for all $g_1,g_2\in G$ (and $\pi(e)=I_m$, where $e\in G$ is the identity element and $I_m$ is the identity matrix). A group element $g\in G$ acts on a vector $v\in V$ by the linear transformation $\pi(g)$, and in this paper we will be dealing with algebraic actions, that is, the entries of the matrix $\pi(g)$ will be rational functions in the entries of the matrix $g$. We will write $g\cdot v=\pi(g)v$. We denote by $\C[V]$ the ring of polynomial functions on $V$.
A polynomial function $f\in\C[V]$ is called invariant if $f(g\cdot v)=f(v)$ for all $v\in V$ and $g\in G$. The induced action on polynomials is $(g\cdot p)(v)=p(g^{-1}\cdot v)$. A set $\{f_i\}_{i\in I}$ of invariant polynomials is called a generating set if any other invariant polynomial can be written as a polynomial in the $f_i$'s. The invariant ring is denoted $\C[V]^G$.
\setcounter{theorem}{2}
\begin{definition}[\textbf{Succinct encoding of generators}] \label{def:succ.enc}
	Fix an action of a group $G$ on a vector space $V = \C^m$. We say that an arithmetic circuit
	$\cC(x_1,\ldots, x_m,y_1,\ldots, y_r)$ \emph{succinctly encodes} the generators of the invariant
	ring if the set of polynomials formed by evaluating the $y$-variables,
	$\{\cC(x_1,\ldots, x_m, \alpha_1,\ldots, \alpha_r)\}_{\alpha_1,\ldots, \alpha_r \in \C}$,
  is a generating set for the invariant ring~$\C[V]^G$.
\end{definition}

\begin{remark}
	The \emph{size} of a succinct encoding as defined above is given by the size of the circuit
	$\cC(x_1, \ldots, x_m, y_1, \ldots, y_r)$, which is measured by the bit complexity of the
	constants used in the computation of $\cC$ as well as the number of gates of the computation
	graph of $\cC$.
  In particular, this means that all constants used in the computation of $\cC$ are rationals.
\end{remark}
